\documentclass[10pt,a4paper]{amsart}
\usepackage{amsmath,amssymb,amsthm,graphicx}
\usepackage{verbatim}
\newenvironment{reference}{\comment}{\endcomment}
%\newenvironment{reference}{}{}
\newenvironment{slogan}{\comment}{\endcomment}
\newenvironment{history}{\comment}{\endcomment}
% For commutative diagrams we use Xy-pic
\usepackage[all]{xy}

% We use 2cell for 2-commutative diagrams.
\xyoption{2cell}
% Theorem environments.
%
\theoremstyle{plain}
\newtheorem{theorem}[subsection]{Theorem}
\newtheorem{proposition}[subsection]{Proposition}
\newtheorem{lemma}[subsection]{Lemma}

\theoremstyle{definition}
\newtheorem{definition}[subsection]{Definition}
\newtheorem{example}[subsection]{Example}
\newtheorem{exercise}[subsection]{Exercise}
\newtheorem{situation}[subsection]{Situation}

\theoremstyle{remark}
\newtheorem{remark}[subsection]{Remark}
\newtheorem{remarks}[subsection]{Remarks}

\numberwithin{equation}{subsection}

% Macros
%
\def\lim{\mathop{\mathrm{lim}}\nolimits}
\def\colim{\mathop{\mathrm{colim}}\nolimits}
\def\Spec{\mathop{\mathrm{Spec}}}
\def\Hom{\mathop{\mathrm{Hom}}\nolimits}
\def\Ext{\mathop{\mathrm{Ext}}\nolimits}
\def\SheafHom{\mathop{\mathcal{H}\!\mathit{om}}\nolimits}
\def\SheafExt{\mathop{\mathcal{E}\!\mathit{xt}}\nolimits}
\def\Sch{\mathit{Sch}}
\def\Mor{\mathop{\mathrm{Mor}}\nolimits}
\def\Ob{\mathop{\mathrm{Ob}}\nolimits}
\def\Sh{\mathop{\mathit{Sh}}\nolimits}
\def\NL{\mathop{N\!L}\nolimits}
\def\CH{\mathop{\mathrm{CH}}\nolimits}
\def\proetale{{pro\text{-}\acute{e}tale}}
\def\etale{{\acute{e}tale}}
\def\QCoh{\mathit{QCoh}}
\def\Ker{\mathop{\mathrm{Ker}}}
\def\Im{\mathop{\mathrm{Im}}}
\def\Coker{\mathop{\mathrm{Coker}}}
\def\Coim{\mathop{\mathrm{Coim}}}

% Boxtimes
%
\DeclareMathSymbol{\boxtimes}{\mathbin}{AMSa}{"02}

%
% Macros for moduli stacks/spaces
%
\def\QCohstack{\mathcal{QC}\!\mathit{oh}}
\def\Cohstack{\mathcal{C}\!\mathit{oh}}
\def\Spacesstack{\mathcal{S}\!\mathit{paces}}
\def\Quotfunctor{\mathrm{Quot}}
\def\Hilbfunctor{\mathrm{Hilb}}
\def\Curvesstack{\mathcal{C}\!\mathit{urves}}
\def\Polarizedstack{\mathcal{P}\!\mathit{olarized}}
\def\Complexesstack{\mathcal{C}\!\mathit{omplexes}}
% \Pic is the operator that assigns to X its picard group, usage \Pic(X)
% \Picardstack_{X/B} denotes the Picard stack of X over B
% \Picardfunctor_{X/B} denotes the Picard functor of X over B
\def\Pic{\mathop{\mathrm{Pic}}\nolimits}
\def\Picardstack{\mathcal{P}\!\mathit{ic}}
\def\Picardfunctor{\mathrm{Pic}}
\def\Deformationcategory{\mathcal{D}\!\mathit{ef}}
\usepackage{cleveref}
\crefname{originallemma}{original lemma}{original lemmas}
\Crefname{originallemma}{Original lemma}{Original lemmas}
\makeatletter
\newlabel{lemma-base-change-gerbe}{{lemma-base-change-gerbe}{source}{}{sourcepointer-1}{}}
\newlabel{lemma-base-change-gerbe@cref}{{[originallemma][0][]lemma-base-change-gerbe}{[1][0][]source}{}{}{}}
\makeatother
\crefname{originallemma}{original lemma}{original lemmas}
\Crefname{originallemma}{Original lemma}{Original lemmas}
\makeatletter
\newlabel{lemma-fppf-open}{{lemma-fppf-open}{source}{}{sourcepointer-2}{}}
\newlabel{lemma-fppf-open@cref}{{[originallemma][0][]lemma-fppf-open}{[1][0][]source}{}{}{}}
\makeatother
\crefname{originallemma}{original lemma}{original lemmas}
\Crefname{originallemma}{Original lemma}{Original lemmas}
\makeatletter
\newlabel{lemma-gerbe-fppf}{{lemma-gerbe-fppf}{source}{}{sourcepointer-3}{}}
\newlabel{lemma-gerbe-fppf@cref}{{[originallemma][0][]lemma-gerbe-fppf}{[1][0][]source}{}{}{}}
\makeatother
\crefname{originalreference}{original reference}{original references}
\Crefname{originalreference}{Original reference}{Original references}
\makeatletter
\newlabel{stacks-lemma-when-gerbe}{{stacks-lemma-when-gerbe}{source}{}{sourcepointer-4}{}}
\newlabel{stacks-lemma-when-gerbe@cref}{{[originalreference][0][]stacks-lemma-when-gerbe}{[1][0][]source}{}{}{}}
\makeatother
\title{A gerbe morphism is a universal homeomorphism}
\author{The Stacks Project Authors}
\date{stacks-project:stacks-morphisms.tex:local-snapshot-20261012; GNU FDL 1.2 or later}
\begin{document}
\maketitle
\noindent\textit{Editorial note.} All delivered mathematical text is the author's own, copied from the pinned Stacks Project chapter source sources/stacks-algebraic-stacks-morphisms/pinned-original-source.tex; no statement or proof step was rewritten, reordered or removed.\par\medskip
\noindent\textit{Editorial note.} Preamble provenance: the Stacks Project keeps one single shared preamble for every chapter; this card takes it directly from the staged sources/stacks-algebraic-stacks-morphisms/preamble.tex as three verbatim spans (lines 8-12, 14-18 and 154-217), so no derived file was written and nothing was added or redefined. The wrapper class and the base packages are supplied by the assembler.\par\medskip
\noindent\textit{Editorial note.} Bibliography provenance: neither the counted statement nor its proof contains a citation (this lane rescans the whole span before writing the card), so the card carries no bibliography and no entry was reconstructed.\par\medskip
\noindent\textit{Editorial note.} Selection provenance: this card was opened under the file-based control contract (CONTROL.md), with the mechanical pre-card gate run on its statement and proof interval before the card was written.\par\medskip
\noindent\textit{Allowed standard prerequisites (editorial).}\begin{itemize}
\item The base change permanence of the gerbe condition. Invoked at the start of the counted proof (source lines 6625-6626). Reduces the assertion to the statement that the map on spaces is a homeomorphism.
\item The two clauses of the gerbe condition. Invoked inside the counted proof (source lines 6631-6633 and 6640-6643). Give fppf local objects for surjectivity and fppf local isomorphisms for injectivity on points.
\item The statements that a gerbe morphism is fppf and that a flat morphism of finite presentation is open on spaces. Invoked at the end of the counted proof (source lines 6646-6647). Supply openness, completing the homeomorphism.
\end{itemize}\medskip
\ifdefined\hypertarget\hypertarget{sourcepointer-1}{}\fi
\noindent\textit{Source reference.} Original lemma lemma-base-change-gerbe: sources/stacks-algebraic-stacks-morphisms/pinned-original-source.tex, source lines 6129--6129. Navigation pointer only: the counted text refers to the Stacks item lemma-base-change-gerbe, which is an existing labeled result of this same chapter and is not part of this card. The pointer carries the target's own label as its number and its definition line as its span; its content is not reproduced here.\par\medskip
\ifdefined\hypertarget\hypertarget{sourcepointer-2}{}\fi
\noindent\textit{Source reference.} Original lemma lemma-fppf-open: sources/stacks-algebraic-stacks-morphisms/pinned-original-source.tex, source lines 5937--5937. Navigation pointer only: the counted text refers to the Stacks item lemma-fppf-open, which is an existing labeled result of this same chapter and is not part of this card. The pointer carries the target's own label as its number and its definition line as its span; its content is not reproduced here.\par\medskip
\ifdefined\hypertarget\hypertarget{sourcepointer-3}{}\fi
\noindent\textit{Source reference.} Original lemma lemma-gerbe-fppf: sources/stacks-algebraic-stacks-morphisms/pinned-original-source.tex, source lines 6287--6287. Navigation pointer only: the counted text refers to the Stacks item lemma-gerbe-fppf, which is an existing labeled result of this same chapter and is not part of this card. The pointer carries the target's own label as its number and its definition line as its span; its content is not reproduced here.\par\medskip
\ifdefined\hypertarget\hypertarget{sourcepointer-4}{}\fi
\noindent\textit{Source reference.} Original reference stacks-lemma-when-gerbe: sources/stacks-algebraic-stacks-morphisms/pinned-original-source.tex, source lines 6111--6111. Navigation pointer only: the counted text refers to the Stacks item stacks-lemma-when-gerbe, which belongs to another chapter of the project and is not part of this card. The pointer carries the target's own label as its number because this chapter's staged snapshot does not contain the other chapter's tag table; its content is not reproduced here.\par\medskip
\begin{lemma}
\label{lemma-gerbe-bijection-points}
Let $f : \mathcal{X} \to \mathcal{Y}$ be a morphism of algebraic stacks.
If $\mathcal{X}$ is a gerbe over $\mathcal{Y}$ then $f$ is a
universal homeomorphism.
\end{lemma}

\begin{proof}
By Lemma \ref{lemma-base-change-gerbe} the assumption on $f$ is
preserved under base change. Hence it suffices to show that the map
$|\mathcal{X}| \to |\mathcal{Y}|$ is a homeomorphism of topological spaces.
Let $k$ be a field and let $y$ be an object
of $\mathcal{Y}$ over $\Spec(k)$. By
Stacks, Lemma \ref{stacks-lemma-when-gerbe} property (2)(a)
there exists an fppf covering $\{T_i \to \Spec(k)\}$
and objects $x_i$ of $\mathcal{X}$ over $T_i$ with $f(x_i) \cong y|_{T_i}$.
Choose an $i$ such that $T_i \not = \emptyset$. Choose a
morphism $\Spec(K) \to T_i$ for some field $K$.
Then $k \subset K$ and $x_i|_K$ is an object of $\mathcal{X}$ lying
over $y|_K$. Thus we see that
$|\mathcal{Y}| \to |\mathcal{X}|$. is surjective. The map
$|\mathcal{Y}| \to |\mathcal{X}|$ is also injective. Namely, if
$x, x'$ are objects of $\mathcal{X}$ over $\Spec(k)$ whose images
$f(x), f(x')$ become isomorphic (over an extension) in $\mathcal{Y}$, then
Stacks, Lemma \ref{stacks-lemma-when-gerbe} property (2)(b)
guarantees the existence of an extension of $k$ over which $x$ and $x'$
become isomorphic (details omitted).
Hence $|\mathcal{X}| \to |\mathcal{Y}|$ is continuous and bijective
and it suffices to show that it is also open.
This follows from
Lemmas \ref{lemma-gerbe-fppf} and \ref{lemma-fppf-open}.
\end{proof}

\end{document}
